

The family $\left(  p_{i,k}\right)  _{\left(  i,k\right)  \in\overline{S}}$ is
summable (by assumption). In other words, for any $m\in\mathbb{N}$, all but
finitely many $\left(  i,k\right)  \in\overline{S}$ satisfy $\left[
x^{m}\right]  p_{i,k}=0$. In other words, for any $m\in\mathbb{N}$, there
exists a finite subset $T_{m}$ of $\overline{S}$ such that
\begin{equation}
\text{all }\left(  i,k\right)  \in\overline{S}\setminus T_{m}\text{ satisfy
}\left[  x^{m}\right]  p_{i,k}=0. \label{pf.prop.fps.prodrule-fin-inf.Tm-def}%
\end{equation}
Consider these finite subsets $T_{m}$.

For each $n\in\mathbb{N}$, we let $T_{n}^{\prime}$ be the subset $T_{0}\cup
T_{1}\cup\cdots\cup T_{n}$ of $\overline{S}$. This subset $T_{n}^{\prime}$ is
finite (since $T_{0},T_{1},\ldots,T_{n}$ are finite).

For each $n\in\mathbb{N}$, we let $I_{n}$ be the subset
\[
\left\{  i\ \mid\ \left(  i,k\right)  \in T_{n}^{\prime}\right\}
\]
of $I$, and we let $K_{n}$ be the set%
\[
\left\{  k\ \mid\ \left(  i,k\right)  \in T_{n}^{\prime}\right\}  .
\]
These two sets $I_{n}$ and $K_{n}$ are finite (since $T_{n}^{\prime}$ is finite).

The definition of $T_{n}^{\prime}$ shows the following:

\begin{statement}
\textit{Claim 3:} Let $n\in\mathbb{N}$ and $\left(  i,k\right)  \in
\overline{S}\setminus T_{n}^{\prime}$. Then, we have%
\[
\left[  x^{m}\right]  p_{i,k}=0\ \ \ \ \ \ \ \ \ \ \text{for each }%
m\in\left\{  0,1,\ldots,n\right\}  .
\]

\end{statement}

[\textit{Proof of Claim 3:} Let $m\in\left\{  0,1,\ldots,n\right\}  $. Then,
$T_{m}\subseteq T_{n}^{\prime}$ (since $T_{n}^{\prime}$ is defined as the
union $T_{0}\cup T_{1}\cup\cdots\cup T_{n}$, whereas $T_{m}$ is one of the
$n+1$ sets appearing in this union). Hence, $T_{n}^{\prime}\supseteq T_{m}$,
so that $\overline{S}\setminus\underbrace{T_{n}^{\prime}}_{\supseteq T_{m}%
}\subseteq\overline{S}\setminus T_{m}$. Now, $\left(  i,k\right)  \in
\overline{S}\setminus T_{n}^{\prime}\subseteq\overline{S}\setminus T_{m}$.
Therefore, (\ref{pf.prop.fps.prodrule-fin-inf.Tm-def}) shows that $\left[
x^{m}\right]  p_{i,k}=0$. This proves Claim 3.]\medskip

The following is easy to see:

\begin{statement}
\textit{Claim 4:} If $\left(  k_{i}\right)  _{i\in I}\in
S_{\operatorname*{fin}}^{I}$, then the family $\left(  p_{i,k_{i}}\right)
_{i\in I}$ is multipliable.
\end{statement}

[\textit{Proof of Claim 4:} Let $\left(  k_{i}\right)  _{i\in I}\in
S_{\operatorname*{fin}}^{I}$. Thus, $\left(  k_{i}\right)  _{i\in I}$ is an
essentially finite family in $S^{I}$. Now, all but finitely many $i\in I$
satisfy $k_{i}=0$ (since $\left(  k_{i}\right)  _{i\in I}$ is essentially
finite) and thus $p_{i,k_{i}}=p_{i,0}=1$ (by
(\ref{eq.prop.fps.prodrule-inf-inf.pi0=1})). Hence, all but finitely many
entries of the family $\left(  p_{i,k_{i}}\right)  _{i\in I}$ equal $1$. Thus,
this family is multipliable (by Proposition \ref{prop.fps.1-mulable}). This
proves Claim 4.]

Claim 4 shows that the product $\prod_{i\in I}p_{i,k_{i}}$ is well-defined
whenever $\left(  k_{i}\right)  _{i\in I}\in S_{\operatorname*{fin}}^{I}$.
Next, we claim the following:

\begin{statement}
\textit{Claim 5:} Let $n\in\mathbb{N}$. Let $\left(  k_{i}\right)  _{i\in
I}\in S_{\operatorname*{fin}}^{I}$. Assume that some $j\in I$ satisfies
$\left(  j,k_{j}\right)  \in\overline{S}\setminus T_{n}^{\prime}$. Then,
$\left[  x^{n}\right]  \left(  \prod_{i\in I}p_{i,k_{i}}\right)  =0$.
\end{statement}

[\textit{Proof of Claim 5:} We have assumed that some $j\in I$ satisfies
$\left(  j,k_{j}\right)  \in\overline{S}\setminus T_{n}^{\prime}$. Consider
this $j$. Then, the product $\prod_{i\in I}p_{i,k_{i}}$ is a multiple of
$p_{j,k_{j}}$ (since $p_{j,k_{j}}$ is one of the factors of this product).
Moreover, we have $\left[  x^{m}\right]  p_{j,k_{j}}=0$ for each $m\in\left\{
0,1,\ldots,n\right\}  $ (by Claim 3, applied to $\left(  i,k\right)  =\left(
j,k_{j}\right)  $). Hence, Lemma \ref{lem.fps.prod.irlv.mul} (applied to
$u=p_{j,k_{j}}$ and $v=\prod_{i\in I}p_{i,k_{i}}$) shows that we have $\left[
x^{m}\right]  \left(  \prod_{i\in I}p_{i,k_{i}}\right)  =0$ for each
$m\in\left\{  0,1,\ldots,n\right\}  $ (since $\prod_{i\in I}p_{i,k_{i}}$ is a
multiple of $p_{j,k_{j}}$). Applying this to $m=n$, we obtain $\left[
x^{n}\right]  \left(  \prod_{i\in I}p_{i,k_{i}}\right)  =0$. This proves Claim 5.]

\begin{statement}
\textit{Claim 6:} Let $n\in\mathbb{N}$. Let $\left(  k_{i}\right)  _{i\in
I}\in S_{\operatorname*{fin}}^{I}\setminus S_{I_{n}}^{I}$. Then, $\left[
x^{n}\right]  \left(  \prod_{i\in I}p_{i,k_{i}}\right)  =0$.
\end{statement}

[\textit{Proof of Claim 6:} We have $\left(  k_{i}\right)  _{i\in I}\in
S_{\operatorname*{fin}}^{I}\setminus S_{I_{n}}^{I}$. In other words, we have
$\left(  k_{i}\right)  _{i\in I}\in S_{\operatorname*{fin}}^{I}$ but $\left(
k_{i}\right)  _{i\in I}\notin S_{I_{n}}^{I}$. From $\left(  k_{i}\right)
_{i\in I}\notin S_{I_{n}}^{I}$, we see that there exists some $j\in I\setminus
I_{n}$ satisfying $k_{j}\neq0$. Consider this $j$. We have $k_{j}\in S_{j}$
(since $\left(  k_{i}\right)  _{i\in I}\in S_{\operatorname*{fin}}%
^{I}\subseteq S^{I}$) and $j\notin I_{n}$ (since $j\in I\setminus I_{n}$). We
have $\left(  j,k_{j}\right)  \in\overline{S}$ (since $k_{j}\in S_{j}$ and
$k_{j}\neq0$) and $\left(  j,k_{j}\right)  \notin T_{n}^{\prime}$ (because if
we had $\left(  j,k_{j}\right)  \in T_{n}^{\prime}$, then we would have $j\in
I_{n}$ by the definition of $I_{n}$; but this would contradict $j\notin I_{n}%
$). Hence, we have $\left(  j,k_{j}\right)  \in\overline{S}\setminus
T_{n}^{\prime}$. Therefore, Claim 5 yields $\left[  x^{n}\right]  \left(
\prod_{i\in I}p_{i,k_{i}}\right)  =0$. This proves Claim 6.]

\begin{statement}
\textit{Claim 7:} The family $\left(  \prod_{i\in I}p_{i,k_{i}}\right)
_{\left(  k_{i}\right)  _{i\in I}\in S_{\operatorname*{fin}}^{I}}$ is summable.
\end{statement}

[\textit{Proof of Claim 7:} Let $n\in\mathbb{N}$. We shall show that all but
finitely many families $\left(  k_{i}\right)  _{i\in I}\in
S_{\operatorname*{fin}}^{I}$ satisfy
\[
\left[  x^{n}\right]  \left(  \prod_{i\in I}p_{i,k_{i}}\right)  =0.
\]


Indeed, let $\left(  k_{i}\right)  _{i\in I}\in S_{\operatorname*{fin}}^{I}$
be a family such that
\begin{equation}
\left[  x^{n}\right]  \left(  \prod_{i\in I}p_{i,k_{i}}\right)  \neq0.
\label{pf.prop.fps.prodrule-fin-inf.c7.pf.ass}%
\end{equation}
We are going to show that $\left(  k_{i}\right)  _{i\in I}$ satisfies the
following two properties:

\begin{itemize}
\item \textit{Property 1:} All $i\in I\setminus I_{n}$ satisfy $k_{i}=0$.

\item \textit{Property 2:} All $i\in I_{n}$ satisfy $k_{i}\in K_{n}%
\cup\left\{  0\right\}  $.
\end{itemize}

\noindent These two properties together will restrict the family $\left(
k_{i}\right)  _{i\in I}$ to finitely many possibilities (since the sets
$I_{n}$ and $K_{n}$ are finite).

If we had $\left(  k_{i}\right)  _{i\in I}\notin S_{I_{n}}^{I}$, then we would
have $\left(  k_{i}\right)  _{i\in I}\in S_{\operatorname*{fin}}^{I}\setminus
S_{I_{n}}^{I}$ (since $\left(  k_{i}\right)  _{i\in I}\in
S_{\operatorname*{fin}}^{I}$) and thus $\left[  x^{n}\right]  \left(
\prod_{i\in I}p_{i,k_{i}}\right)  =0$ (by Claim 6), which would contradict
(\ref{pf.prop.fps.prodrule-fin-inf.c7.pf.ass}). Hence, we cannot have $\left(
k_{i}\right)  _{i\in I}\notin S_{I_{n}}^{I}$. Thus, we have $\left(
k_{i}\right)  _{i\in I}\in S_{I_{n}}^{I}$. In other words, all $i\in
I\setminus I_{n}$ satisfy $k_{i}=0$. This proves Property 1.

If there was some $j\in I_{n}$ that satisfies $k_{j}\notin K_{n}\cup\left\{
0\right\}  $, then this $k_{j}$ would satisfy $\left(  j,k_{j}\right)
\in\overline{S}\setminus T_{n}^{\prime}$\ \ \ \ \footnote{\textit{Proof.} Let
$j\in I_{n}$ be such that $k_{j}\notin K_{n}\cup\left\{  0\right\}  $. We must
show that $\left(  j,k_{j}\right)  \in\overline{S}\setminus T_{n}^{\prime}$.
\par
Indeed, we have $k_{j}\notin K_{n}\cup\left\{  0\right\}  $; thus,
$k_{j}\notin K_{n}$ and $k_{j}\neq0$. However, $j\in I_{n}\subseteq I$ and
$k_{j}\in S_{j}$ (since $\left(  k_{i}\right)  _{i\in I}\in
S_{\operatorname*{fin}}^{I}\subseteq S^{I}$) and therefore $\left(
j,k_{j}\right)  \in\overline{S}$ (by the definition of $\overline{S}$, since
$k_{j}\neq0$). If we had $\left(  j,k_{j}\right)  \in T_{n}^{\prime}$, then we
would have $k_{j}\in K_{n}$ (by the definition of $K_{n}$), which would
contradict $k_{j}\notin K_{n}$. Thus, we have $\left(  j,k_{j}\right)  \notin
T_{n}^{\prime}$. Combining this with $\left(  j,k_{j}\right)  \in\overline{S}%
$, we obtain $\left(  j,k_{j}\right)  \in\overline{S}\setminus T_{n}^{\prime}%
$. This completes our proof.}, and therefore we would have $\left[
x^{n}\right]  \left(  \prod_{i\in I}p_{i,k_{i}}\right)  =0$ (by Claim 5, since
$j\in I_{n}\subseteq I$); but this would contradict
(\ref{pf.prop.fps.prodrule-fin-inf.c7.pf.ass}). Hence, there exists no $j\in
I_{n}$ that satisfies $k_{j}\notin K_{n}\cup\left\{  0\right\}  $. In other
words, all $i\in I_{n}$ satisfy $k_{i}\in K_{n}\cup\left\{  0\right\}  $. This
proves Property 2.

Now, we have shown that our family $\left(  k_{i}\right)  _{i\in I}$ satisfies
Property 1 and Property 2.

Forget that we fixed $\left(  k_{i}\right)  _{i\in I}$. We thus have shown
that any family $\left(  k_{i}\right)  _{i\in I}\in S_{\operatorname*{fin}%
}^{I}$ that satisfies $\left[  x^{n}\right]  \left(  \prod_{i\in I}p_{i,k_{i}%
}\right)  \neq0$ must satisfy Property 1 and Property 2. In other words, any
such family must belong to the set of all families $\left(  k_{i}\right)
_{i\in I}\in S_{\operatorname*{fin}}^{I}$ that satisfy Property 1 and Property
2. However, the latter set is finite\footnote{\textit{Proof.} We must show
that Property 1 and Property 2 leave only finitely many options for the family
$\left(  k_{i}\right)  _{i\in I}$. Indeed, Property 1 shows that all entries
$k_{i}$ with $i\in I\setminus I_{n}$ are uniquely determined; meanwhile,
Property 2 ensures that the remaining entries (of which there are only
finitely many, since the set $I_{n}$ is finite) must belong to the finite set
$K_{n}\cup\left\{  0\right\}  $ (this set is finite, since $K_{n}$ is finite).
Therefore, a family $\left(  k_{i}\right)  _{i\in I}$ that satisfies Property
1 and Property 2 is uniquely determined by finitely many of its entries
(namely, by its entries $k_{i}$ with $i\in I_{n}$), and there are finitely
many choices for each of them (since they must belong to the finite set
$K_{n}\cup\left\{  0\right\}  $). Hence, there are only finitely many such
families (namely, at most $\left\vert K_{n}\cup\left\{  0\right\}  \right\vert
^{\left\vert I_{n}\right\vert }$ many options). In other words, the set of all
families $\left(  k_{i}\right)  _{i\in I}\in S_{\operatorname*{fin}}^{I}$ that
satisfy Property 1 and Property 2 is finite.}. Hence, there are only finitely
many families $\left(  k_{i}\right)  _{i\in I}\in S_{\operatorname*{fin}}^{I}$
that satisfy $\left[  x^{n}\right]  \left(  \prod_{i\in I}p_{i,k_{i}}\right)
\neq0$ (since we have shown that any such family must belong to the finite set
of all families $\left(  k_{i}\right)  _{i\in I}\in S_{\operatorname*{fin}%
}^{I}$ that satisfy Property 1 and Property 2). In other words, all but
finitely many families $\left(  k_{i}\right)  _{i\in I}\in
S_{\operatorname*{fin}}^{I}$ satisfy $\left[  x^{n}\right]  \left(
\prod_{i\in I}p_{i,k_{i}}\right)  =0$.

Forget that we fixed $n$. We thus have shown that for each $n\in\mathbb{N}$,
all but finitely many families $\left(  k_{i}\right)  _{i\in I}\in
S_{\operatorname*{fin}}^{I}$ satisfy $\left[  x^{n}\right]  \left(
\prod_{i\in I}p_{i,k_{i}}\right)  =0$. In other words, the family $\left(
\prod_{i\in I}p_{i,k_{i}}\right)  _{\left(  k_{i}\right)  _{i\in I}\in
S_{\operatorname*{fin}}^{I}}$ is summable. This proves Claim 7.]

Claim 7 shows that the sum $\sum_{\left(  k_{i}\right)  _{i\in I}\in
S_{\operatorname*{fin}}^{I}}\ \ \prod_{i\in I}p_{i,k_{i}}$ is well-defined. We
shall next focus on proving (\ref{eq.prop.fps.prodrule-inf-inf.eq}).

\begin{statement}
\textit{Claim 8:} Let $n\in\mathbb{N}$. Then,
\[
\left[  x^{n}\right]  \left(  \sum_{\left(  k_{i}\right)  _{i\in I}\in
S_{\operatorname*{fin}}^{I}}\ \ \prod_{i\in I}p_{i,k_{i}}\right)  =\left[
x^{n}\right]  \left(  \sum_{\left(  k_{i}\right)  _{i\in I_{n}}\in S^{I_{n}}%
}\ \ \prod_{i\in I_{n}}p_{i,k_{i}}\right)  .
\]

\end{statement}

[\textit{Proof of Claim 8:} The set $I_{n}$ is finite (as we know). Thus,
$S_{I_{n}}^{I}\subseteq S_{\operatorname*{fin}}^{I}$%
\ \ \ \ \footnote{\textit{Proof.} Let $\left(  k_{i}\right)  _{i\in I}\in
S_{I_{n}}^{I}$. Thus, $\left(  k_{i}\right)  _{i\in I}$ is a family in $S^{I}$
that satisfies $k_{i}=0$ for all $i\in I\setminus I_{n}$ (by the definition of
$S_{I_{n}}^{I}$). However, $I_{n}$ is finite. Thus, $k_{i}=0$ for all but
finitely many $i\in I$ (since $k_{i}=0$ for all $i\in I\setminus I_{n}$). In
other words, the family $\left(  k_{i}\right)  _{i\in I}$ is essentially
finite. In other words, $\left(  k_{i}\right)  _{i\in I}\in
S_{\operatorname*{fin}}^{I}$ (by the definition of $S_{\operatorname*{fin}%
}^{I}$).
\par
Forget that we fixed $\left(  k_{i}\right)  _{i\in I}$. We thus have shown
that $\left(  k_{i}\right)  _{i\in I}\in S_{\operatorname*{fin}}^{I}$ for each
$\left(  k_{i}\right)  _{i\in I}\in S_{I_{n}}^{I}$. In other words, $S_{I_{n}%
}^{I}\subseteq S_{\operatorname*{fin}}^{I}$.}. Hence, the set
$S_{\operatorname*{fin}}^{I}$ is the union of its two disjoint subsets
$S_{I_{n}}^{I}$ and $S_{\operatorname*{fin}}^{I}\setminus S_{I_{n}}^{I}$.

Furthermore, for each $\left(  k_{i}\right)  _{i\in I_{n}}\in S_{I_{n}}^{I}$,
we have%
\[
k_{i}=0\ \ \ \ \ \ \ \ \ \ \text{for all }i\in I\setminus I_{n}%
\]
(by the definition of $S_{I_{n}}^{I}$) and therefore%
\[
p_{i,k_{i}}=p_{i,0}=1\ \ \ \ \ \ \ \ \ \ \text{for all }i\in I\setminus I_{n}%
\]
(by (\ref{eq.prop.fps.prodrule-inf-inf.pi0=1})) and thus%
\[
\prod_{i\in I\setminus I_{n}}\underbrace{p_{i,k_{i}}}_{=1}=\prod_{i\in
I\setminus I_{n}}1=1
\]
and therefore%
\begin{align}
\prod_{i\in I}p_{i,k_{i}}  &  =\left(  \prod_{i\in I_{n}}p_{i,k_{i}}\right)
\underbrace{\left(  \prod_{i\in I\setminus I_{n}}p_{i,k_{i}}\right)  }%
_{=1}\nonumber\\
&  \ \ \ \ \ \ \ \ \ \ \ \ \ \ \ \ \ \ \ \ \left(
\begin{array}
[c]{c}%
\text{here, we have split the product into two}\\
\text{parts, since the set }I_{n}\text{ is a subset of }I
\end{array}
\right) \nonumber\\
&  =\prod_{i\in I_{n}}p_{i,k_{i}}.
\label{pf.prop.fps.prodrule-fin-inf.c8.pf.1}%
\end{align}


Now,%
\begin{align*}
&  \left[  x^{n}\right]  \left(  \sum_{\left(  k_{i}\right)  _{i\in I}\in
S_{\operatorname*{fin}}^{I}}\ \ \prod_{i\in I}p_{i,k_{i}}\right) \\
&  =\sum_{\left(  k_{i}\right)  _{i\in I}\in S_{\operatorname*{fin}}^{I}%
}\left[  x^{n}\right]  \left(  \prod_{i\in I}p_{i,k_{i}}\right) \\
&  =\sum_{\left(  k_{i}\right)  _{i\in I}\in S_{I_{n}}^{I}}\left[
x^{n}\right]  \underbrace{\left(  \prod_{i\in I}p_{i,k_{i}}\right)
}_{\substack{=\prod_{i\in I_{n}}p_{i,k_{i}}\\\text{(by
(\ref{pf.prop.fps.prodrule-fin-inf.c8.pf.1}))}}}+\sum_{\left(  k_{i}\right)
_{i\in I}\in S_{\operatorname*{fin}}^{I}\setminus S_{I_{n}}^{I}}%
\underbrace{\left[  x^{n}\right]  \left(  \prod_{i\in I}p_{i,k_{i}}\right)
}_{\substack{=0\\\text{(by Claim 6)}}}\\
&  \ \ \ \ \ \ \ \ \ \ \ \ \ \ \ \ \ \ \ \ \left(
\begin{array}
[c]{c}%
\text{here, we have split the sum, since the set }S_{\operatorname*{fin}}%
^{I}\text{ is}\\
\text{the union of its two disjoint subsets }S_{I_{n}}^{I}\text{ and
}S_{\operatorname*{fin}}^{I}\setminus S_{I_{n}}^{I}%
\end{array}
\right) \\
&  =\sum_{\left(  k_{i}\right)  _{i\in I}\in S_{I_{n}}^{I}}\left[
x^{n}\right]  \left(  \prod_{i\in I_{n}}p_{i,k_{i}}\right)  +\underbrace{\sum
_{\left(  k_{i}\right)  _{i\in I}\in S_{\operatorname*{fin}}^{I}\setminus
S_{I_{n}}^{I}}0}_{=0}\\
&  =\sum_{\left(  k_{i}\right)  _{i\in I}\in S_{I_{n}}^{I}}\left[
x^{n}\right]  \left(  \prod_{i\in I_{n}}p_{i,k_{i}}\right)  =\sum_{\left(
k_{i}\right)  _{i\in I_{n}}\in S^{I_{n}}}\left[  x^{n}\right]  \left(
\prod_{i\in I_{n}}p_{i,k_{i}}\right) \\
&  \ \ \ \ \ \ \ \ \ \ \ \ \ \ \ \ \ \ \ \ \left(
\begin{array}
[c]{c}%
\text{here, we have substituted }\left(  k_{i}\right)  _{i\in I_{n}}\text{ for
}\left(  k_{i}\right)  _{i\in I_{n}}\text{ in the}\\
\text{sum, since the map }S_{I_{n}}^{I}\rightarrow S^{I_{n}},\ \left(
k_{i}\right)  _{i\in I}\mapsto\left(  k_{i}\right)  _{i\in I_{n}}\\
\text{is a bijection (indeed, this map is the map}\\
\text{we have called }\operatorname*{reduce}\nolimits_{I_{n}}\text{)}%
\end{array}
\right) \\
&  =\left[  x^{n}\right]  \left(  \sum_{\left(  k_{i}\right)  _{i\in I_{n}}\in
S^{I_{n}}}\ \ \prod_{i\in I_{n}}p_{i,k_{i}}\right)  .
\end{align*}
This proves Claim 8.]

\begin{statement}
\textit{Claim 9:} Let $n\in\mathbb{N}$. Let $i\in I\setminus I_{n}$. Then,
\[
\left[  x^{m}\right]  \left(  \sum_{k\in S_{i}\setminus\left\{  0\right\}
}p_{i,k}\right)  =0\ \ \ \ \ \ \ \ \ \ \text{for each }m\in\left\{
0,1,\ldots,n\right\}  .
\]

\end{statement}

[\textit{Proof of Claim 9:} We have $i\in I\setminus I_{n}$. In other words,
$i\in I$ and $i\notin I_{n}$.

Let $m\in\left\{  0,1,\ldots,n\right\}  $.

Let $k\in S_{i}\setminus\left\{  0\right\}  $. Thus, by the definition of
$\overline{S}$, we have $\left(  i,k\right)  \in\overline{S}$ (since $k\in
S_{i}\setminus\left\{  0\right\}  $ entails $k\in S_{i}$ and $k\neq0$). On the
other hand, $\left(  i,k\right)  \notin T_{n}^{\prime}$ (since otherwise, we
would have $\left(  i,k\right)  \in T_{n}^{\prime}$ and thus $i\in I_{n}$ (by
the definition of $I_{n}$), which would contradict $i\notin I_{n}$). Combining
$\left(  i,k\right)  \in\overline{S}$ with $\left(  i,k\right)  \notin
T_{n}^{\prime}$, we obtain $\left(  i,k\right)  \in\overline{S}\setminus
T_{n}^{\prime}$. Therefore, Claim 3 yields $\left[  x^{m}\right]  p_{i,k}=0$.

Now, forget that we fixed $k$. We thus have shown that
\begin{equation}
\left[  x^{m}\right]  p_{i,k}=0\ \ \ \ \ \ \ \ \ \ \text{for each }k\in
S_{i}\setminus\left\{  0\right\}  .
\label{pf.prop.fps.prodrule-fin-inf.c9.pf.1}%
\end{equation}
Now,%
\[
\left[  x^{m}\right]  \left(  \sum_{k\in S_{i}\setminus\left\{  0\right\}
}p_{i,k}\right)  =\sum_{k\in S_{i}\setminus\left\{  0\right\}  }%
\underbrace{\left[  x^{m}\right]  p_{i,k}}_{\substack{=0\\\text{(by
(\ref{pf.prop.fps.prodrule-fin-inf.c9.pf.1}))}}}=\sum_{k\in S_{i}%
\setminus\left\{  0\right\}  }0=0.
\]
This proves Claim 9.]

\begin{statement}
\textit{Claim 10:} Let $n\in\mathbb{N}$. Then,%
\[
\left[  x^{n}\right]  \left(  \prod_{i\in I}\ \ \sum_{k\in S_{i}}%
p_{i,k}\right)  =\left[  x^{n}\right]  \left(  \prod_{i\in I_{n}}%
\ \ \sum_{k\in S_{i}}p_{i,k}\right)  .
\]

\end{statement}

[\textit{Proof of Claim 10:} The set $I_{n}$ is a subset of $I$. Hence, the
set $I$ is the union of its two disjoint subsets $I_{n}$ and $I\setminus
I_{n}$. Thus, we can split the product $\prod_{i\in I}\ \ \sum_{k\in S_{i}%
}p_{i,k}$ as follows:\footnote{Here, we are tacitly using the fact that any
subfamily of the family $\left(  \sum_{k\in S_{i}}p_{i,k}\right)  _{i\in I}$
is multipliable. But this can be proved in the exact same way as we proved
Claim 2.}%
\begin{align}
\prod_{i\in I}\ \ \sum_{k\in S_{i}}p_{i,k}  &  =\left(  \prod_{i\in I_{n}%
}\ \ \sum_{k\in S_{i}}p_{i,k}\right)  \left(  \prod_{i\in I\setminus I_{n}%
}\ \ \underbrace{\sum_{k\in S_{i}}p_{i,k}}_{\substack{=1+\sum_{k\in
S_{i}\setminus\left\{  0\right\}  }p_{i,k}\\\text{(by
(\ref{pf.prop.fps.prodrule-fin-inf.sum-no-1}))}}}\right) \nonumber\\
&  =\left(  \prod_{i\in I_{n}}\ \ \sum_{k\in S_{i}}p_{i,k}\right)  \left(
\prod_{i\in I\setminus I_{n}}\left(  1+\sum_{k\in S_{i}\setminus\left\{
0\right\}  }p_{i,k}\right)  \right)  .
\label{pf.prop.fps.prodrule-fin-inf.c10.pf.1}%
\end{align}


However, the family $\left(  \sum_{k\in S_{i}\setminus\left\{  0\right\}
}p_{i,k}\right)  _{i\in I}$ is summable (as we have seen in the proof of Claim
2). Hence, its subfamily $\left(  \sum_{k\in S_{i}\setminus\left\{  0\right\}
}p_{i,k}\right)  _{i\in I\setminus I_{n}}$ is summable as well (since a
subfamily of a summable family is always summable). Moreover, Claim 9 shows
that each $i\in I\setminus I_{n}$ satisfies $\left[  x^{m}\right]  \left(
\sum_{k\in S_{i}\setminus\left\{  0\right\}  }p_{i,k}\right)  =0$ for each
$m\in\left\{  0,1,\ldots,n\right\}  $. Hence, Lemma
\ref{lem.fps.prod.irlv.inf} (applied to $a=\prod_{i\in I_{n}}\ \ \sum_{k\in
S_{i}}p_{i,k}$ and $J=I\setminus I_{n}$ and $f_{i}=\sum_{k\in S_{i}%
\setminus\left\{  0\right\}  }p_{i,k}$) yields that%
\begin{align*}
&  \left[  x^{m}\right]  \left(  \left(  \prod_{i\in I_{n}}\ \ \sum_{k\in
S_{i}}p_{i,k}\right)  \left(  \prod_{i\in I\setminus I_{n}}\left(
1+\sum_{k\in S_{i}\setminus\left\{  0\right\}  }p_{i,k}\right)  \right)
\right) \\
&  =\left[  x^{m}\right]  \left(  \prod_{i\in I_{n}}\ \ \sum_{k\in S_{i}%
}p_{i,k}\right)  \ \ \ \ \ \ \ \ \ \ \text{for each }m\in\left\{
0,1,\ldots,n\right\}  .
\end{align*}
Applying this to $m=n$, we find%
\[
\left[  x^{n}\right]  \left(  \left(  \prod_{i\in I_{n}}\ \ \sum_{k\in S_{i}%
}p_{i,k}\right)  \left(  \prod_{i\in I\setminus I_{n}}\left(  1+\sum_{k\in
S_{i}\setminus\left\{  0\right\}  }p_{i,k}\right)  \right)  \right)  =\left[
x^{n}\right]  \left(  \prod_{i\in I_{n}}\ \ \sum_{k\in S_{i}}p_{i,k}\right)
.
\]
In view of (\ref{pf.prop.fps.prodrule-fin-inf.c10.pf.1}), this rewrites as%
\[
\left[  x^{n}\right]  \left(  \prod_{i\in I}\ \ \sum_{k\in S_{i}}%
p_{i,k}\right)  =\left[  x^{n}\right]  \left(  \prod_{i\in I_{n}}%
\ \ \sum_{k\in S_{i}}p_{i,k}\right)  .
\]
This proves Claim 10.]

\begin{statement}
\textit{Claim 11:} Let $i\in\mathbb{N}$. Then,%
\[
\prod_{i\in I_{n}}\ \ \sum_{k\in S_{i}}p_{i,k}=\sum_{\left(  k_{i}\right)
_{i\in I_{n}}\in S^{I_{n}}}\ \ \prod_{i\in I_{n}}p_{i,k_{i}}.
\]

\end{statement}

[\textit{Proof of Claim 11:} The set $I_{n}$ is finite. For any $i\in I_{n}$,
the family $\left(  p_{i,k}\right)  _{k\in S_{i}}$ is summable (by Claim 1).
Hence, Proposition \ref{prop.fps.prodrule-fin-infJ} (applied to $N=I_{n}$)
yields
\[
\prod_{i\in I_{n}}\ \ \sum_{k\in S_{i}}p_{i,k}=\sum_{\left(  k_{i}\right)
_{i\in I_{n}}\in\prod_{i\in I_{n}}S_{i}}\ \ \prod_{i\in I_{n}}p_{i,k_{i}}%
=\sum_{\left(  k_{i}\right)  _{i\in I_{n}}\in S^{I_{n}}}\ \ \prod_{i\in I_{n}%
}p_{i,k_{i}}%
\]
(since $\prod_{i\in I_{n}}S_{i}=S^{I_{n}}$). This proves Claim 11.]\medskip

Now, for each $n\in\mathbb{N}$, we have%
\begin{align*}
&  \left[  x^{n}\right]  \left(  \prod_{i\in I}\ \ \sum_{k\in S_{i}}%
p_{i,k}\right) \\
&  =\left[  x^{n}\right]  \left(  \prod_{i\in I_{n}}\ \ \sum_{k\in S_{i}%
}p_{i,k}\right)  \ \ \ \ \ \ \ \ \ \ \left(  \text{by Claim 10}\right) \\
&  =\left[  x^{n}\right]  \left(  \sum_{\left(  k_{i}\right)  _{i\in I_{n}}\in
S^{I_{n}}}\ \ \prod_{i\in I_{n}}p_{i,k_{i}}\right) \\
&  \ \ \ \ \ \ \ \ \ \ \ \ \ \ \ \ \ \ \ \ \left(  \text{since Claim 11 yields
}\prod_{i\in I_{n}}\ \ \sum_{k\in S_{i}}p_{i,k}=\sum_{\left(  k_{i}\right)
_{i\in I_{n}}\in S^{I_{n}}}\ \ \prod_{i\in I_{n}}p_{i,k_{i}}\right) \\
&  =\left[  x^{n}\right]  \left(  \sum_{\left(  k_{i}\right)  _{i\in I}\in
S_{\operatorname*{fin}}^{I}}\ \ \prod_{i\in I}p_{i,k_{i}}\right)
\ \ \ \ \ \ \ \ \ \ \left(  \text{by Claim 8}\right)  .
\end{align*}
That is, any coefficient of the FPS $\prod_{i\in I}\ \ \sum_{k\in S_{i}%
}p_{i,k}$ equals the corresponding coefficient of $\sum_{\left(  k_{i}\right)
_{i\in I}\in S_{\operatorname*{fin}}^{I}}\ \ \prod_{i\in I}p_{i,k_{i}}$.
Hence,%
\[
\prod_{i\in I}\ \ \sum_{k\in S_{i}}p_{i,k}=\sum_{\left(  k_{i}\right)  _{i\in
I}\in S_{\operatorname*{fin}}^{I}}\ \ \prod_{i\in I}p_{i,k_{i}}=\sum
_{\substack{\left(  k_{i}\right)  _{i\in I}\in\prod_{i\in I}S_{i}\\\text{is
essentially finite}}}\ \ \prod_{i\in I}p_{i,k_{i}}%
\]
(since $S_{\operatorname*{fin}}^{I}$ is the set of all essentially finite
families $\left(  k_{i}\right)  _{i\in I}\in\prod_{i\in I}S_{i}$). In
particular, the family $\left(  \prod_{i\in I}p_{i,k_{i}}\right)  _{\left(
k_{i}\right)  _{i\in I}\in\prod_{i\in I}S_{i}\text{ is essentially finite}}$
is summable. This proves Proposition \ref{prop.fps.prodrule-inf-inf}.
\end{proof}
\end{fineprint}

\begin{fineprint}
Our proof of Proposition \ref{prop.fps.subs.rule-infprod} will use the
\textbf{finite} analogue of Proposition \ref{prop.fps.subs.rule-infprod},
which is easy:

\begin{lemma}
\label{lem.fps.subs.rule-infprod-fin}Let $I$ be a \textbf{finite} set. If
$\left(  f_{i}\right)  _{i\in I}\in K\left[  \left[  x\right]  \right]  ^{I}$
is a family of FPSs, and if $g\in K\left[  \left[  x\right]  \right]  $ is an
FPS satisfying $\left[  x^{0}\right]  g=0$, then $\left(  \prod_{i\in I}%
f_{i}\right)  \circ g=\prod_{i\in I}\left(  f_{i}\circ g\right)  $.
\end{lemma}

\begin{proof}
[Proof of Lemma \ref{lem.fps.subs.rule-infprod-fin}.]This follows by a
straightforward induction on $\left\vert I\right\vert $. (The base case is the
case when $\left\vert I\right\vert =0$, and relies on the fact that
$\underline{1}\circ g=\underline{1}$ for any $g\in K\left[  \left[  x\right]
\right]  $. The induction step relies on Proposition \ref{prop.fps.subs.rules}
\textbf{(b)}. The details are left to the reader, who must have seen dozens of
such proofs by now.)
\end{proof}

\begin{proof}
[Detailed proof of Proposition \ref{prop.fps.subs.rule-infprod}.]Let $\left(
f_{i}\right)  _{i\in I}\in K\left[  \left[  x\right]  \right]  ^{I}$ be a
multipliable family of FPSs. Let $g\in K\left[  \left[  x\right]  \right]  $
be an FPS satisfying $\left[  x^{0}\right]  g=0$.

We shall first show the following auxiliary claim:

\begin{statement}
\textit{Claim 1:} Let $M$ be an $x^{n}$-approximator for $\left(
f_{i}\right)  _{i\in I}$. Then, the set $M$ determines the $x^{n}$-coefficient
in the product of $\left(  f_{i}\circ g\right)  _{i\in I}$.
\end{statement}

[\textit{Proof of Claim 1:} The set $M$ is an $x^{n}$-approximator for
$\left(  f_{i}\right)  _{i\in I}$. In other words, $M$ is a finite subset of
$I$ that determines the first $n+1$ coefficients in the product of $\left(
f_{i}\right)  _{i\in I}$ (by the definition of \textquotedblleft$x^{n}%
$-approximator\textquotedblright).

Let $J$ be a finite subset of $I$ satisfying $M\subseteq J\subseteq I$. Then,
Lemma \ref{lem.fps.subs.rule-infprod-fin} (applied to $J$ instead of $I$)
yields%
\begin{equation}
\left(  \prod_{i\in J}f_{i}\right)  \circ g=\prod_{i\in J}\left(  f_{i}\circ
g\right)  . \label{pf.prop.fps.subs.rule-infprod.3}%
\end{equation}
Also, Lemma \ref{lem.fps.subs.rule-infprod-fin} (applied to $M$ instead of
$I$) yields%
\begin{equation}
\left(  \prod_{i\in M}f_{i}\right)  \circ g=\prod_{i\in M}\left(  f_{i}\circ
g\right)  . \label{pf.prop.fps.subs.rule-infprod.4}%
\end{equation}
However, Proposition \ref{prop.fps.infprod-approx-xneq} \textbf{(a)} (applied
to $\mathbf{a}_{i}=f_{i}$) yields%
\[
\prod_{i\in J}f_{i}\overset{x^{n}}{\equiv}\prod_{i\in M}f_{i}.
\]
We also have $g\overset{x^{n}}{\equiv}g$ (since the relation $\overset{x^{n}%
}{\equiv}$ is an equivalence relation). Hence, Proposition
\ref{prop.fps.xneq.comp} (applied to $a=\prod_{i\in J}f_{i}$ and
$b=\prod_{i\in M}f_{i}$ and $c=g$ and $d=g$) yields%
\[
\left(  \prod_{i\in J}f_{i}\right)  \circ g\overset{x^{n}}{\equiv}\left(
\prod_{i\in M}f_{i}\right)  \circ g.
\]
In view of (\ref{pf.prop.fps.subs.rule-infprod.3}) and
(\ref{pf.prop.fps.subs.rule-infprod.4}), this rewrites as%
\[
\prod_{i\in J}\left(  f_{i}\circ g\right)  \overset{x^{n}}{\equiv}\prod_{i\in
M}\left(  f_{i}\circ g\right)  .
\]
In other words,
\[
\text{each }m\in\left\{  0,1,\ldots,n\right\}  \text{ satisfies }\left[
x^{m}\right]  \left(  \prod_{i\in J}\left(  f_{i}\circ g\right)  \right)
=\left[  x^{m}\right]  \left(  \prod_{i\in M}\left(  f_{i}\circ g\right)
\right)
\]
(by the definition of the relation $\overset{x^{n}}{\equiv}$). Applying this
to $m=n$, we obtain%
\[
\left[  x^{n}\right]  \left(  \prod_{i\in J}\left(  f_{i}\circ g\right)
\right)  =\left[  x^{n}\right]  \left(  \prod_{i\in M}\left(  f_{i}\circ
g\right)  \right)  .
\]


Forget that we fixed $J$. We thus have shown that every finite subset $J$ of
$I$ satisfying $M\subseteq J\subseteq I$ satisfies%
\[
\left[  x^{n}\right]  \left(  \prod_{i\in J}\left(  f_{i}\circ g\right)
\right)  =\left[  x^{n}\right]  \left(  \prod_{i\in M}\left(  f_{i}\circ
g\right)  \right)  .
\]
In other words, the set $M$ determines the $x^{n}$-coefficient in the product
of $\left(  f_{i}\circ g\right)  _{i\in I}$ (by the definition of what it
means to \textquotedblleft determine the $x^{n}$-coefficient in the product of
$\left(  f_{i}\circ g\right)  _{i\in I}$\textquotedblright). This proves Claim
1.] \medskip

Now, let $n\in\mathbb{N}$. Lemma \ref{lem.fps.mulable.approx} (applied to
$\mathbf{a}_{i}=f_{i}$) shows that there exists an $x^{n}$-approximator for
$\left(  f_{i}\right)  _{i\in I}$. Consider this $x^{n}$-approximator for
$\left(  f_{i}\right)  _{i\in I}$, and denote it by $M$. Thus, $M$ is an
$x^{n}$-approximator for $\left(  f_{i}\right)  _{i\in I}$; in other words,
$M$ is a finite subset of $I$ that determines the first $n+1$ coefficients in
the product of $\left(  f_{i}\right)  _{i\in I}$ (by the definition of
\textquotedblleft$x^{n}$-approximator\textquotedblright). Claim 1 shows that
the set $M$ determines the $x^{n}$-coefficient in the product of $\left(
f_{i}\circ g\right)  _{i\in I}$. Hence, there is a finite subset of $I$ that
determines the $x^{n}$-coefficient in the product of $\left(  f_{i}\circ
g\right)  _{i\in I}$ (namely, $M$). In other words, the $x^{n}$-coefficient in
the product of $\left(  f_{i}\circ g\right)  _{i\in I}$ is finitely determined.

Forget that we fixed $n$. We thus have shown that for each $n\in\mathbb{N}$,
the $x^{n}$-coefficient in the product of $\left(  f_{i}\circ g\right)  _{i\in
I}$ is finitely determined. In other words, each coefficient in the product of
$\left(  f_{i}\circ g\right)  _{i\in I}$ is finitely determined. In other
words, the family $\left(  f_{i}\circ g\right)  _{i\in I}$ is multipliable (by
the definition of \textquotedblleft multipliable\textquotedblright). \medskip

It remains to prove that $\left(  \prod_{i\in I}f_{i}\right)  \circ
g=\prod_{i\in I}\left(  f_{i}\circ g\right)  $.

In order to do so, we again fix $n\in\mathbb{N}$. Lemma
\ref{lem.fps.mulable.approx} (applied to $\mathbf{a}_{i}=f_{i}$) shows that
there exists an $x^{n}$-approximator for $\left(  f_{i}\right)  _{i\in I}$.
Consider this $x^{n}$-approximator for $\left(  f_{i}\right)  _{i\in I}$, and
denote it by $M$. Thus, $M$ is an $x^{n}$-approximator for $\left(
f_{i}\right)  _{i\in I}$; in other words, $M$ is a finite subset of $I$ that
determines the first $n+1$ coefficients in the product of $\left(
f_{i}\right)  _{i\in I}$ (by the definition of \textquotedblleft$x^{n}%
$-approximator\textquotedblright). Moreover, Proposition
\ref{prop.fps.infprod-approx-xneq} \textbf{(b)} (applied to $\mathbf{a}%
_{i}=f_{i}$) yields%
\begin{equation}
\prod_{i\in I}f_{i}\overset{x^{n}}{\equiv}\prod_{i\in M}f_{i}.
\label{pf.prop.fps.subs.rule-infprod.1}%
\end{equation}
We also have $g\overset{x^{n}}{\equiv}g$ (since the relation $\overset{x^{n}%
}{\equiv}$ is an equivalence relation). Hence, Proposition
\ref{prop.fps.xneq.comp} (applied to $a=\prod_{i\in I}f_{i}$ and
$b=\prod_{i\in M}f_{i}$ and $c=g$ and $d=g$) yields%
\begin{equation}
\left(  \prod_{i\in I}f_{i}\right)  \circ g\overset{x^{n}}{\equiv}\left(
\prod_{i\in M}f_{i}\right)  \circ g. \label{pf.prop.fps.subs.rule-infprod.2}%
\end{equation}
However, Lemma \ref{lem.fps.subs.rule-infprod-fin} (applied to $M$ instead of
$I$) yields%
\[
\left(  \prod_{i\in M}f_{i}\right)  \circ g=\prod_{i\in M}\left(  f_{i}\circ
g\right)  .
\]
In view of this, we can rewrite (\ref{pf.prop.fps.subs.rule-infprod.2}) as
\[
\left(  \prod_{i\in I}f_{i}\right)  \circ g\overset{x^{n}}{\equiv}\prod_{i\in
M}\left(  f_{i}\circ g\right)  .
\]
In other words,
\[
\text{each }m\in\left\{  0,1,\ldots,n\right\}  \text{ satisfies }\left[
x^{m}\right]  \left(  \left(  \prod_{i\in I}f_{i}\right)  \circ g\right)
=\left[  x^{m}\right]  \left(  \prod_{i\in M}\left(  f_{i}\circ g\right)
\right)
\]
(by the definition of the relation $\overset{x^{n}}{\equiv}$). Applying this
to $m=n$, we obtain%
\begin{equation}
\left[  x^{n}\right]  \left(  \left(  \prod_{i\in I}f_{i}\right)  \circ
g\right)  =\left[  x^{n}\right]  \left(  \prod_{i\in M}\left(  f_{i}\circ
g\right)  \right)  . \label{pf.prop.fps.subs.rule-infprod.6}%
\end{equation}


However, Claim 1 shows that the set $M$ determines the $x^{n}$-coefficient in
the product of $\left(  f_{i}\circ g\right)  _{i\in I}$. Hence, the definition
of the infinite product $\prod_{i\in I}\left(  f_{i}\circ g\right)  $
(specifically, Definition \ref{def.fps.multipliable} \textbf{(b)}) yields%
\[
\left[  x^{n}\right]  \left(  \prod_{i\in I}\left(  f_{i}\circ g\right)
\right)  =\left[  x^{n}\right]  \left(  \prod_{i\in M}\left(  f_{i}\circ
g\right)  \right)  .
\]
Comparing this with (\ref{pf.prop.fps.subs.rule-infprod.6}), we obtain%
\[
\left[  x^{n}\right]  \left(  \left(  \prod_{i\in I}f_{i}\right)  \circ
g\right)  =\left[  x^{n}\right]  \left(  \prod_{i\in I}\left(  f_{i}\circ
g\right)  \right)  .
\]


Forget that we fixed $n$. We thus have shown that each $n\in\mathbb{N}$
satisfies%
\[
\left[  x^{n}\right]  \left(  \left(  \prod_{i\in I}f_{i}\right)  \circ
g\right)  =\left[  x^{n}\right]  \left(  \prod_{i\in I}\left(  f_{i}\circ
g\right)  \right)  .
\]
In other words, the FPSs $\left(  \prod_{i\in I}f_{i}\right)  \circ g$ and
$\prod_{i\in I}\left(  f_{i}\circ g\right)  $ agree in all their coefficients.
Hence, $\left(  \prod_{i\in I}f_{i}\right)  \circ g=\prod_{i\in I}\left(
f_{i}\circ g\right)  $. This completes our proof of Proposition
\ref{prop.fps.subs.rule-infprod}.
\end{proof}
\end{fineprint}
